\section{Conclusion}
\label{sec:conclusion}

% Target: ~0.25 two-column page (~0.5 single-column equivalent)
% Condense report sections/conclusion.tex.

% [Paragraph 1: Summary of what was done.]
%   This paper investigated the feasibility of integrating time protection
%   mechanisms into a mainstream operating system by extending FreeRTOS-MPU
%   with a domain-level scheduler. The scheduler partitions execution into
%   fixed time slices across security domains and employs the RISC-V
%   temporal fence (fence.t) instruction on every domain switch to clear
%   shared microarchitectural state, thereby closing scheduler-based and
%   shared-state timing channels.

% [Paragraph 2: Key results.]
%   Evaluation on QEMU system emulation, using instruction-based cycle
%   counting, demonstrates constant-time domain switching with $\pm 1$ cycle
%   variance across two workloads. The multi-domain demo confirms that the
%   scheduler preserves intra-domain priority semantics while enforcing
%   strict temporal isolation at domain boundaries.

% [Paragraph 3: Limitations and closing.]
%   The evaluation is limited to instruction-level determinism; validation
%   of microarchitectural isolation requires real hardware or cycle-accurate
%   simulation. Further work is needed to audit system-call timing, partition
%   higher-level caches, and explore lighter-weight alternatives that balance
%   protection and throughput. Nevertheless, this work demonstrates that
%   time protection is not exclusive to clean-slate microkernels: it can be
%   integrated into an existing, widely deployed RTOS, bringing meaningful
%   timing-channel mitigation to mainstream embedded systems.
